\chapter{Límite central, órbitas de Kepler y holonomía de Sommerfeld}
\label{chap:limite-central-kepler-sommerfeld}
\index{Kepler, órbita de}\index{Sommerfeld, cuantización de}
\index{Levi-Civita, cubierta de}\index{Möbius, holonomía de}

Este capítulo distingue cuatro estratos que la terminología física puede
superponer indebidamente: la cancelación exacta de flujos interiores en un
grafo; la realización continua de un campo central; la dinámica kepleriana
que resulta del potencial inverso; y la condición semiclásica de fase sobre
una banda con holonomía. La primera es combinatoria. Las tres restantes
requieren, sucesivamente, una carta métrica, un lagrangiano y una línea de
fase semiclásica.

\section{Gauss discreto y descenso inverso al área}
\label{sec:gauss-discreto-limite-central}

Sea \(G=(V,E)\) un grafo finito no orientado, sea
\(\vec E:=\{(x,y),(y,x):\{x,y\}\in E\}\) su conjunto de arcos y sea
\(F:\vec E\to\mathbb R\) una \(1\)-co\-cadena antisimétrica,
\(F(y,x)=-F(x,y)\). Definimos
\[
 \operatorname{div}F(x)
 :=\sum_{y:\{x,y\}\in E}F(x,y).
\]
Para \(S\subseteq V\), el flujo
\(\operatorname{Flux}_{\partial S}F\) es la suma orientada sobre las aristas
con un extremo en \(S\) y el otro en \(V\setminus S\).

\begin{theorem}[Teorema de Gauss en grafos]
\label{thm:gauss-grafo-hmt-md}
Para todo \(S\subseteq V\),
\[
 \sum_{x\in S}\operatorname{div}F(x)
 =\operatorname{Flux}_{\partial S}F.
 \tag{K1}\label{eq:gauss-grafo-hmt-md}
\]
\end{theorem}

\begin{proof}
Cada arista con ambos extremos en \(S\) aparece dos veces con orientaciones
opuestas y se cancela. Sólo permanecen las aristas que atraviesan la cadena
de borde.
\end{proof}

En la red cúbica, para \(R\in\mathbb Z_{\geq0}\),
\[
 C_R:=\{x\in\mathbb Z^3:\|x\|_\infty\leq R\},
\]
el número de enlaces salientes es
\[
 |\partial C_R|=6(2R+1)^2.
\]
Si una fuente posee flujo total \(Q\) y se impone uniformidad media sobre el
cascarón, el flujo por enlace vale
\[
 E_R=\frac{Q}{6(2R+1)^2}.
\]
Al definir \(r_{\rm eff}\) mediante
\[
 4\pi r_{\rm eff}^2=6(2R+1)^2,
\]
se obtiene la identidad exacta
\[
 E_R=\frac{Q}{4\pi r_{\rm eff}^2}.
 \tag{K2}\label{eq:inverso-area-cascaron}
\]
El paso desde \eqref{eq:gauss-grafo-hmt-md} a
\eqref{eq:inverso-area-cascaron} usa la hipótesis adicional de isotropía
media. Identificar \(Q\) con una fuente gravitatoria o eléctrica requiere,
además, el transductor dimensional correspondiente.

\section{Lagrangiano central y cónicas de Kepler}
\label{sec:lagrangiano-central-kepler}

Sea \(\mu>0\) una masa reducida y
\(\mathcal K>0\) un acoplamiento del tipo producto de energía y longitud,
\([\mathcal K]=\mathsf E\,\mathsf L\). En la hoja
euclídea continua consideremos
\[
 L_{\rm Kep}(\mathbf r,\dot{\mathbf r})
 =\frac{\mu}{2}\|\dot{\mathbf r}\|^2+\frac{\mathcal K}{r},
 \qquad
 H_{\rm Kep}(\mathbf r,\mathbf p)
 =\frac{\|\mathbf p\|^2}{2\mu}-\frac{\mathcal K}{r},
 \tag{K3}\label{eq:lagrangiano-hamiltoniano-kepler}
\]
donde \(r=\|\mathbf r\|>0\) y
\(\mathbf p=\mu\dot{\mathbf r}\). La ecuación de Euler--Lagrange es
\[
 \mu\ddot{\mathbf r}
 =-\frac{\mathcal K}{r^3}\mathbf r.
 \tag{K4}\label{eq:ecuacion-central-kepler}
\]
El momento angular específico
\(\mathbf h=\mathbf r\times\dot{\mathbf r}\) es constante; si
\(\mathbf h\ne0\), la trayectoria permanece en el plano ortogonal a
\(\mathbf h\).

\begin{proposition}[Cónica kepleriana]
\label{prop:conica-kepleriana-hmt-md}
Sea \(h=\|\mathbf h\|\). Toda solución no radial de
\eqref{eq:ecuacion-central-kepler} satisface
\[
 r(\phi)
 =\frac{p}{1+e\cos(\phi-\phi_0)},
 \qquad
 p=\frac{\mu h^2}{\mathcal K}.
 \tag{K5}\label{eq:conica-kepleriana}
\]
Para energía negativa, \(0\leq e<1\) y la órbita es elíptica.
\end{proposition}

\begin{proof}
Con \(u=1/r\), la ecuación de Binet adopta la forma
\[
 u''(\phi)+u(\phi)=\frac{\mathcal K}{\mu h^2}.
\]
Su solución general es
\[
 u=\frac{\mathcal K}{\mu h^2}
 \bigl(1+e\cos(\phi-\phi_0)\bigr).
\]
Invertir \(u\) da \eqref{eq:conica-kepleriana}; la condición de ligadura
equivale a \(e<1\).
\end{proof}

La velocidad areolar es
\[
 \frac{\mathrm dA}{\mathrm dt}=\frac h2.
 \tag{K6}\label{eq:segunda-ley-kepler}
\]
Para una elipse de semieje mayor \(a\), integrar
\eqref{eq:segunda-ley-kepler} durante un periodo produce
\[
 A_{\rm ell}=\pi ab,
 \qquad
 b=a\sqrt{1-e^2},
 \qquad
 p=a(1-e^2)=\frac{\mu h^2}{\mathcal K}.
\]
Por tanto,
\(T=2A_{\rm ell}/h\), y al eliminar \(b\), \(p\) y \(h\) se obtiene
\[
 T^2=\frac{4\pi^2\mu}{\mathcal K}\,a^3.
 \tag{K7}\label{eq:tercera-ley-kepler}
\]
En el problema gravitatorio de prueba,
\(\mathcal K=G M\mu\), y
\eqref{eq:tercera-ley-kepler} se reduce a
\(T^2=4\pi^2a^3/(GM)\) \cite{Kepler1609,Newton1687}.

El teorema de Gauss del grafo es exacto. Bajo isotropía y distancia
declaradas, el flujo produce el campo inverso al área; el límite continuo
transporta después ese campo al lagrangiano central y a las órbitas
keplerianas escritas arriba.

\section{Cubierta de Levi--Civita y banda de Möbius}
\label{sec:levi-civita-mobius-kepler}

En el plano complejo sin el origen, la aplicación de Levi--Civita
\[
 q=z^2
 \tag{K8}\label{eq:cubierta-levi-civita}
\]
es una doble cubierta. Una vuelta de \(q\) alrededor del origen intercambia
los dos levantamientos \(z\) y \(-z\); hacen falta dos vueltas para cerrar el
levantamiento. Es
un testigo exacto de memoria \(\mathbb Z_2\). No es, sin embargo, una banda
de Möbius: la regularización completa del problema de Kepler exige también
reparametrizar el tiempo y transformar los momentos.

Sea ahora \(\gamma:S^1\hookrightarrow\mathbb R^2\) un embedding cuya imagen
es una órbita elíptica, y fijemos \(\varepsilon>0\). Definimos
\[
 M_\gamma
 :=([0,2\pi]\times[-\varepsilon,\varepsilon])/
 \bigl((0,u)\sim(2\pi,-u)\bigr).
 \tag{K9}\label{eq:banda-mobius-orbita}
\]

\begin{theorem}[Periodo orientado de la banda orbital]
\label{thm:periodo-orientado-banda-orbital}
\(M_\gamma\) es una banda de Möbius cuyo núcleo es \(\gamma\). Una sección
de la línea de signo posee holonomía
\[
 \mathcal H(\gamma)=-I,
 \qquad
 \mathcal H(\gamma^2)=I.
 \tag{K10}\label{eq:holonomia-banda-orbital}
\]
Por tanto, su retorno como sección orientada requiere dos vueltas.
\end{theorem}

\begin{proof}
La identificación de las fibras extremas mediante \(u\mapsto-u\) es la
construcción estándar de la banda de Möbius. La holonomía de una vuelta es
la involución de la fibra; su cuadrado es la identidad.
\end{proof}

La cubierta \eqref{eq:cubierta-levi-civita} y la banda
\eqref{eq:banda-mobius-orbita} comparten una monodromía de orden dos, pero no
son objetos isomorfos. Esta distinción impide deducir una orientación de
partícula--antipartícula a partir de la regularización de Levi--Civita sin un
mapa adicional.

\section{Condición semiclásica de Sommerfeld con holonomía}
\label{sec:sommerfeld-holonomia-banda}

Sea
\[
 J_\gamma:=\oint_\gamma\mathbf p\cdot\mathrm d\mathbf q
\]
la acción de un ciclo y supóngase que el transporte semiclásico aporta la
fase \(\exp(\mathrm iJ_\gamma/\hbar_{\rm int})\). Si la línea real asociada
posee holonomía \(\varepsilon_\gamma\in\{+1,-1\}\), la condición de cierre de
la sección complejificada es
\[
 \varepsilon_\gamma
 \exp\!\left(\frac{\mathrm iJ_\gamma}{\hbar_{\rm int}}\right)=1.
 \tag{K11}\label{eq:cierre-fase-holonomia-sommerfeld}
\]

\begin{proposition}[Desplazamiento de media unidad por holonomía de signo]
\label{prop:desplazamiento-sommerfeld-holonomia}
Bajo las hipótesis anteriores,
\[
 \begin{array}{lll}
 \varepsilon_\gamma=+1
 &\Longrightarrow&
 J_\gamma=2\pi\hbar_{\rm int}n,\\[1mm]
 \varepsilon_\gamma=-1
 &\Longrightarrow&
 J_\gamma=2\pi\hbar_{\rm int}\left(n+\dfrac12\right),
 \end{array}
 \qquad n\in\mathbb Z.
 \tag{K12}\label{eq:cuantizacion-sommerfeld-holonomia}
\]
\end{proposition}

\begin{proof}
La \cref{eq:cierre-fase-holonomia-sommerfeld} exige una fase igual a uno en
el sector trivial e igual a menos uno en el sector de signo. Sus argumentos
difieren, respectivamente, de un múltiplo entero o semientero de \(2\pi\).
\end{proof}

La proposición es una consecuencia exacta de la fase y de la holonomía
declaradas; no selecciona por sí sola el Hamiltoniano de Coulomb, un número
cuántico físico ni un espectro atómico \cite{Sommerfeld1916}. En una
cuantización EBK deben añadirse los índices de Maslov de los ciclos
correspondientes, sin contar dos veces el desplazamiento topológico.

Relativamente al lector de orientación, las dos hojas locales pueden
representar las rutas \(\gamma\) y \(\gamma^\vee\) y enlazarse con las
dos orientaciones del entrelazamiento bicapa definido en
\cref{def:entropia-bicapa-rev6}. Esta correspondencia no identifica
esas rutas con funciones de Green avanzadas o retardadas ni afirma que la
entropía cause la cónica kepleriana: la órbita procede de
\eqref{eq:ecuacion-central-kepler}; la entropía mide, después de una
completación de Hilbert, la correlación entre hojas.

\begin{criterion}[Límite Kepler--Sommerfeld]
\label{crit:limite-kepler-sommerfeld}
La realización queda refutada o indebidamente tipada si se obtiene una ley
inversa al cuadrado sin declarar isotropía; si
\(\mathcal K\) o \(G\) se introducen como escalares sin recta dimensional;
si la cubierta \(q=z^2\) se identifica con una banda de Möbius; si
\eqref{eq:cuantizacion-sommerfeld-holonomia} se publica como espectro exacto
sin Hamiltoniano, ciclos e índices de Maslov; o si se atribuye
entrelazamiento a una mera inversión determinista de ruta.
\end{criterion}
